% Hafta 2 — Olay incelemelerinin ortak iskeleti
% Derleme: py -3.12 tools/latex/derle.py docs/week-2/assets/h02-11-olay-iskeleti.tex
\documentclass[tikz,border=8pt]{standalone}
\usepackage{cen429-cizim}
\begin{document}
\begin{tikzpicture}[node distance=10mm and 9mm]
  \node[baslik] (b) at (0,0) {Olay incelemelerinin ortak iskeleti};
  \node[altbaslik] at (b.south west) {Morris'ten WannaCry'a kadar aynı zincir};

  \node[kotukutu=38mm, minimum height=24mm, below=13mm of b.south west, anchor=north west, xshift=6mm] (p1)
    {\kb{Programlama hatası}\\taşma, doğrulanmamış girdi,\\eksik yama};
  \node[kotukutu=38mm, minimum height=24mm, right=of p1] (p2)
    {\kb{Zararlı yazılım girer}\\solucan, fidye, arka kapı};
  \node[kotukutu=38mm, minimum height=24mm, right=of p2] (p3)
    {\kb{Zarar}\\veri kaybı, üretim durması,\\itibar};
  \node[iyikutu=38mm, minimum height=24mm, right=of p3] (p4)
    {\kb{Önlem}\\güvenli kodlama + yama +\\bütünlük izleme};

  \draw[kotuok] (p1) -- (p2);
  \draw[kotuok] (p2) -- (p3);
  \draw[ok, draw=soluk] (p3) -- (p4);

  \node[dip, text width=185mm, below=10mm of p2.south, anchor=north]
    {Her olay incelemesinde bu dört kutuyu doldurun; ders zincirin neresinde kırıldığıdır.};
\end{tikzpicture}
\end{document}
